\documentclass[12pt]{article}
\usepackage[utf8]{inputenc}
\usepackage[T1]{fontenc}
\usepackage[brazil]{babel}
\usepackage[a4paper,margin=2.5cm]{geometry}
\usepackage{enumitem}
\usepackage{amsmath,amssymb}
\setlist[itemize]{noitemsep,topsep=4pt}
\title{Sprint 3 - S3\_12\\Erros Numéricos}
\author{Matheus Moreli Durão Silva\\Matrícula 12621CBS217\\Cibersegurança - Bacharelado - Noturno}
\date{2026}
\begin{document}
\maketitle

\section*{Entregável}
Experimentos com ponto flutuante e estabilidade numérica. Este documento registra o objetivo do card, a solução desenvolvida e a forma de conferência. Os arquivos de código e as evidências citadas permanecem na mesma pasta deste documento.

\section*{O que foi estudado e aplicado}
\begin{itemize}
\item Erro absoluto e relativo
\item Truncamento e arredondamento
\item Tolerância numérica
\item Número de condição e matriz quase singular
\item Comparação solve e inv
\end{itemize}

\section*{Como conferir}
O avaliador pode abrir os arquivos desta pasta, executar os exemplos quando houver código e comparar a saída com a explicação apresentada. Nos cards ligados ao Inventário de Ativos, a solução foi relacionada ao Trabalho 2 sem modificar requisitos já atendidos do CRUD.

\section*{Representação e aproximação}
Em binary64, $0,1+0,2$ é representado aproximadamente por $0,30000000000000004$. O erro absoluto contra $0,3$ é cerca de $5,55\cdot10^{-17}$; o relativo é cerca de $1,85\cdot10^{-16}$. Para $6,8085106$, truncar a duas casas dá $6,80$ e arredondar dá $6,81$; os erros absolutos são aproximadamente $0,0085106$ e $0,0014894$.
\section*{Residual no risco}
Para $I-A=\left[\begin{smallmatrix}1&-0,2\\-0,3&1\end{smallmatrix}\right]$ e $p=(8,4)^T$, $\det(I-A)=0,94$ e a solução é $(440/47,320/47)^T$. O script calcula $e=(I-A)r-p$ e sua norma máxima, compara \texttt{solve} com inversão explícita e mede o número de condição. Residual pequeno confirma equações computadas, mas matriz mal condicionada pode amplificar erro de entrada.
\section*{Matriz quase singular da aula}
Para $A_2=\left[\begin{smallmatrix}0&0,1&0,9\\0,1&0&0,9\\0,1&0,9&0\end{smallmatrix}\right]$, cada linha soma 1 e $I-A_2$ é singular em teoria. Em ponto flutuante, seu determinante pode sair próximo de zero sem ser exatamente zero. Por isso o script verifica $\kappa_2(I-A_2)$ com \texttt{np.linalg.cond}: acima de $1/\varepsilon_{\mathrm{máq}}$, rejeita o sistema; acima de $10^8$, avisa sobre perda de precisão. O programa principal usa $\kappa_\infty$, uma norma diferente, e também resolve o sistema sem calcular a inversa para obter o risco.
\section*{Conclusão}
O artefato foi organizado para demonstrar tanto o conceito estudado quanto sua aplicação prática. As evidências de execução só são consideradas válidas quando produzidas pela ferramenta correspondente.
\end{document}
